% =============================================================================
% Results section excerpt -- assembled from the manuscript draft, with all
% figure-vs-supplement edits applied (EC50, Emax, k_immune, transfer rates).
% This is NOT a full paper: it holds only the Results section so it can be
% compiled standalone for review. Merge \section{Results} below into the
% main manuscript .tex file, and make sure the main file's own preamble
% loads amsmath, graphicx, hyperref, and cleveref (all used here).
% =============================================================================
\documentclass[11pt]{article}

\usepackage[margin=1in]{geometry}
\usepackage{amsmath,amssymb}
\usepackage{graphicx}
\usepackage[colorlinks=true,linkcolor=blue,citecolor=blue]{hyperref}
\usepackage[capitalize]{cleveref}
\usepackage{float}

\title{Results (excerpt)}
\author{}
\date{}

\begin{document}
\maketitle

\section{Results}

Prior to initiation of Vancomycin, bacterial populations established a
quasi-steady-state equilibrium across both compartments. The
vancomycin-sensitive \textit{S. aureus} strain in the blood compartment
exhibited a growth rate of $\rho_s = 0.61\text{ h}^{-1}$ while the resistant
strain growth was $\rho_r = 0.55\text{ h}^{-1}$. In the bone reservoir, local
tissue conditions reduced bacterial replication rates to $\rho_s^r =
0.175\text{ h}^{-1}$ for the sensitive strain and $\rho_r^r = 0.1765\text{
h}^{-1}$ for the resistant strain. Asymmetric exchange between compartments
leads to seeding of the blood compartment from the reservoir.

The clinical simulation captures the transition from osteomyelitis to
secondary bacteremia and subsequent infective endocarditis. Vancomycin
administration (1000~mg every 8 hours), models clinically recommended
steady-state peak and trough plasma concentrations. Vancomycin therapy drove
rapid killing of susceptible \textit{S. aureus} subpopulations ($S_b$ and
$S_{\text{res}}$). However, vancomycin exerted no growth inhibition on
resistant strains ($R_b$ and $R_{\text{res}}$; MIC = 32~mg/dL). Following
laboratory confirmation of vancomycin-resistant MRSA, therapy was switched to
Linezolid (600~mg every 12 hours), yielding recommended steady-state serum
concentrations. Linezolid administration drove a reduction in bacterial
counts across both compartments. However, due to reduced local immune
efficacy, bacterial clearance from the reservoir lagged significantly behind
vascular clearance.

Monte Carlo simulation of linezolid potency ($EC_{50,L}$), maximum killing
capacity ($E_{max,L}$), host immune clearance ($k_{immune}$) and transfer
rates between compartments each varied a single mechanistic parameter around
its calibrated baseline while holding all other parameters fixed. All
samples were propagated through the full 81-day dual-reservoir model to
quantify how uncertainty in each parameter translates into heterogeneity in
clinical outcome.

% -----------------------------------------------------------------------
% EC50,L
% -----------------------------------------------------------------------
Sampling $EC_{50,L}$ from a log-normal distribution with mean
$1.0~\mathrm{mg/L}$ and $\mathrm{CV}=1.117$ produced escape of both the
blood ($R_b$) and reservoir ($R_{res}$) resistant subpopulations in 34.7\%
of simulations. The $R_b$ and $R_{res}$ escaped in the identical set of runs
in every case, consistent with a single shared escape boundary at
$EC_{50,L} \approx 0.93~\mathrm{mg/L}$ (\Cref{fig:ec50_terminal_burden}),
located by a companion deterministic threshold sweep and reproduced
independently by the continuous Monte Carlo switch curve (Supplementary
Figs.~S1, S2). Because roughly a third of the sampled $EC_{50,L}$ values
(median $0.70$, 95th percentile $2.81~\mathrm{mg/L}$) exceeded this
threshold, reduced linezolid potency was sufficient on its own to convert a
clearing infection into a persistent one, independent of variation in any
other model parameter.

\begin{figure}[H]
    \centering
    \includegraphics[width=1.0\linewidth]{figures/ec50_terminal_burden_vs_ec50.png}
    \caption{Terminal $R_{res}$ burden versus $EC_{50,L}$: open circles
    pinned to the LOD for runs that clear, filled circles at their actual
    final count for runs that escape, split by the bisected switch at
    $0.932~\mathrm{mg/L}$.}
    \label{fig:ec50_terminal_burden}
\end{figure}

% -----------------------------------------------------------------------
% Emax,L
% -----------------------------------------------------------------------
Sampling $E_{max,L}$ from a log-normal distribution (mean
$0.8~\mathrm{h}^{-1}$ and $\mathrm{CV} = 0.30$) produced an escape rate of
56.5\% (565/1000) for both $R_b$ and $R_{res}$. Because the bisected
threshold ($0.8161~\mathrm{h}^{-1}$) sits just \emph{above} the sampled
mean, just over half of simulated patients drew an $E_{max,L}$ value too low
to clear the infection, with the population split occurring at the same
threshold as the deterministic sweep (Supplementary Fig.~S12).

The deterministic threshold sweep located a shared escape boundary at
$E_{max,L} \approx 0.8161~\mathrm{h}^{-1}$ (Supplementary Fig.~S11): both
the blood ($R_b$) and reservoir ($R_{res}$) compartments escaped below this
value and cleared above it, crossing the limit of detection (LOD) at the
same $E_{max,L}$, consistent with the single escape boundary observed across
all four sensitivity analyses. \Cref{fig:emax_terminal_burden} shows the
switch point where the terminal $R_{res}$ burden moves from the LOD floor to
its plateau at the bisected threshold.

\begin{figure}[H]
    \centering
    \includegraphics[width=1.0\linewidth]{figures/emax_terminal_burden_vs_emax.png}
    \caption{Terminal $R_{res}$ burden versus $E_{max,L}$: open circles
    pinned to the LOD for runs that clear, filled circles at their actual
    final count for runs that escape, split by the bisected switch at
    $0.8161~\mathrm{h}^{-1}$.}
    \label{fig:emax_terminal_burden}
\end{figure}

% -----------------------------------------------------------------------
% k_immune
% -----------------------------------------------------------------------
Sampling $k_{immune}$ from a log-normal distribution with mean
$0.17~\mathrm{h}^{-1}$ and $\mathrm{CV} = 0.25$ produced a relapse rate of
35.7\% (357/1000): both compartments remained above the LOD at the end of
follow-up in these runs, while the remaining 64.3\% cleared entirely, with
the population split occurring exactly at the same threshold as the
deterministic sweep (Supplementary Fig.~S10). Because this threshold sits
below the sampled mean, most simulated patients draw a $k_{immune}$ value
sufficient to clear the resistant subpopulations -- but with $\mathrm{CV} =
0.25$, over a third draw a value too low to do so, indicating that
resistant-strain persistence in this model remains meaningfully sensitive to
inter-patient variability in host immune clearance even though the majority
of the population resolves the infection.

The deterministic threshold sweep located a shared escape boundary at
$k_{immune} \approx 0.15~\mathrm{h}^{-1}$ (Supplementary Fig.~S9): both the
blood ($R_b$) and reservoir ($R_{res}$) compartments escaped below this
value and cleared above it, crossing the limit of detection (LOD) at the
same $k_{immune}$ -- a shared threshold rather than two independent ones,
consistent with the single escape boundary observed across all four
sensitivity analyses. The model's own default, $k_{immune} =
0.12~\mathrm{h}^{-1}$, falls well within the escape zone.

\Cref{fig:mc_immune_response} shows the resulting full time-course kinetics
as the median and 5th--95th percentile envelope across all 1000 sampled
trajectories. Both sensitive subpopulations ($S_b$, $S_{res}$) rise to their
pretreatment plateau and clear below the LOD during the vancomycin window,
regardless of the sampled $k_{immune}$ value. $R_b$ shows a modest early
peak that declines through the remainder of the pretreatment period, then
spikes sharply during the vancomycin-to-linezolid transition before clearing
for the full 42-day linezolid course; $R_{res}$ follows the same
transition-period spike but only partially clears, settling at a
reduced-but-detectable plateau (on the order of $10^2~\mathrm{CFU/mL}$) for
the remainder of therapy. Both resistant compartments relapse to their
pretreatment/carrying-capacity levels within days of linezolid
discontinuation, consistent with the reservoir acting as a persistent source
that reseeds blood once treatment pressure is removed.

\begin{figure}[H]
    \centering
    \includegraphics[width=1.0\linewidth]{figures/mc_immune_response.png}
    \caption{Monte Carlo sweep of $k_{immune}$ ($n = 1000$): full
    time-course kinetics for $S_b$, $R_b$, $S_{res}$, and $R_{res}$, shown
    as individual trajectories (light), median (solid), and 5th--95th
    percentile envelope (shaded).}
    \label{fig:mc_immune_response}
\end{figure}

% -----------------------------------------------------------------------
% Transfer rates
% -----------------------------------------------------------------------
Full time-course kinetics for both sweeps, shown as the median and
5th--95th percentile envelope across all 193 grid points, are given in
\Cref{fig:mc_frb_sweep_kinetics,fig:mc_fbr_sweep_kinetics}. Because 160 of
the 193 grid points are concentrated in the $[10^{-3},
10^{-2}]\,\mathrm{h}^{-1}$ band -- a band where, as detailed below, both
compartments are already established -- these median trajectories are
dominated by near-baseline behavior rather than by the low-rate regime
detailed below: $S_b$ and $S_{res}$ clear below the LOD during the combined
vancomycin/linezolid course, $R_b$ transiently exceeds the LOD during the
vancomycin-to-linezolid transition before also clearing, and $R_{res}$
persists at a reduced but detectable level throughout linezolid therapy
before regrowing to its reservoir carrying capacity after treatment
discontinuation, reseeding a detectable $R_b$ relapse by the end of
follow-up.

\begin{figure}[H]
    \centering
    \includegraphics[width=1.0\linewidth]{figures/mc_frb_sweep_kinetics.png}
    \caption{Full time-course kinetics, Sweep 1: $f_{r \to b}$ (reservoir
    $\to$ blood) varied across the log-spaced grid, $f_{b \to r}$ held at
    baseline ($1\times10^{-5}\,\mathrm{h}^{-1}$). Lines show the median and
    5th--95th percentile envelope across all 193 grid points.}
    \label{fig:mc_frb_sweep_kinetics}
\end{figure}

\begin{figure}[H]
    \centering
    \includegraphics[width=1.0\linewidth]{figures/mc_fbr_sweep_kinetics.png}
    \caption{Full time-course kinetics, Sweep 2: $f_{b \to r}$ (blood $\to$
    reservoir) varied across the log-spaced grid, $f_{r \to b}$ held at
    baseline ($5\times10^{-5}\,\mathrm{h}^{-1}$). Lines show the median and
    5th--95th percentile envelope across all 193 grid points.}
    \label{fig:mc_fbr_sweep_kinetics}
\end{figure}

For $f_{r \to b}$ (Sweep 1), peak $S_b$, peak $R_b$, and final $R_b$
remained at or near zero for $f_{r \to b} \lesssim
3\times10^{-5}\,\mathrm{h}^{-1}$: blood is seeded exclusively through
reservoir-to-blood transfer (blood compartments start sterile), so
vanishingly small $f_{r \to b}$ leaves bacteremia unestablished for the
entire simulated course. All three then switch sharply to their full
baseline levels within roughly half a log-decade of the model's own
baseline value ($5\times10^{-5}\,\mathrm{h}^{-1}$; Supplementary Fig.~S4;
sampling grid in Supplementary Fig.~S8). Peak $S_{res}$ is the only one of
the four outcomes that is genuinely flat across the full six-decade range,
varying by less than 10\% from $10^{-8}$ to $10^{-2}\,\mathrm{h}^{-1}$. At
the opposite end of the range, final $R_b$ undergoes a second, independent
transition: it drops back to zero for $f_{r \to b} \gtrsim
5\times10^{-3}\,\mathrm{h}^{-1}$ even as peak $R_b$ continues to rise,
indicating that $R_b$ still establishes and peaks at these rates but no
longer relapses by the end of follow-up (Supplementary Fig.~S5).

For $f_{b \to r}$ (Sweep 2), peak $S_b$, peak $R_b$, and final $R_b$ are
flat to within 1\% across roughly the lower four and a half of the six
log-decades sampled ($10^{-8}$ to $10^{-3.5}\,\mathrm{h}^{-1}$;
Supplementary Fig.~S6), consistent with $f_{b \to r}$ acting only as an
efflux from blood rather than a seeding route for either compartment. Peak
$S_{res}$, however, is not flat over this same upper range: it rises from
its baseline value of $\sim$4560~CFU/mL to a maximum of $\sim$6270~CFU/mL
($+$37\%) near $f_{b \to r} \approx 3\times10^{-3}\,\mathrm{h}^{-1}$, before
collapsing back toward baseline as peak $S_b$ simultaneously collapses
toward zero, both by $f_{b \to r} = 1\times10^{-2}\,\mathrm{h}^{-1}$
(Supplementary Fig.~S7).

\end{document}
